\section*{\texorpdfstring{\S\CurrentExerciseGroup}{第{\CurrentExerciseGroup}节}}
\phantomsection
\addcontentsline{toc}{section}{第{\CurrentExerciseGroup}节习题}

\begin{enumerate}
\def\labelenumi{\arabic{enumi})}
\tightlist
\item
  \begin{enumerate}
  \def\labelenumii{\alph{enumii})}
  \tightlist
  \item
    设 \(\mathbf{K}\) 是 Hausdorff 局部凸空间 \(\mathbf{E}\) 的一个紧凸子集，\(\mu\) 是 \(\mathbf{K}\) 上质量为 1 的一个正测度，\(\mathbf{b}_{\mu}\) 是它的重心。试证：对每个正函数 \(f\) （\(\mathbf{K}\) 上的、凹的下半连续者），有 \(f(\mathbf{b}_{\mu}) \geqslant \int^{*} f d\mu\) 。（注意 \(f\) 有界（TVS, II, §2 习题 32），因而 \(-f\) 是 \(\mu\) -可积的。）由此推出：若 \(f\) 是 \(\mathbf{K}\) 上既凸又凹的下半连续（或上半连续）函数，则 \(\int f d\mu = f(\mathbf{b}_{\mu})\) 。
  \end{enumerate}
\end{enumerate}

\begin{enumerate}
\def\labelenumi{\alph{enumi})}
\setcounter{enumi}{1}
\tightlist
\item
  设 I 是 R 中的区间 \([0,1]\) ，K 是 \(\mathcal{M}_{+}(I)\) 中由总质量为 1 的正测度所组成的子集，它对淡拓扑是凸的且紧的，并设 \(j: x \mapsto \varepsilon_{x}\) 是 I 到 K 中的典范单射，它是 I 到 K 的一个子空间上的同胚。对每个测度 \(\nu \in K\) ，令 \(g(\nu) = \sum_{x \in I} \nu(\{x\})\) ；这是一个
\end{enumerate}

在 K 中既凸又凹的函数；此外，若 \(g_{n}(\nu)\) 是对 I 的至多含 n 个元素的所有有限子集 A 取的诸 \(\nu(A)\) 的上确界，则 \(g_{n}\) 在 K 上上半连续，且 \(g(\nu) = \lim_{n \to \infty} g_{n}(\nu)\) 对每个 \(\nu \in K\) 成立，因此 g 对 K 上每个测度 \(\lambda\) 都是 \(\lambda\) -可积的。设 \(\mu\) 是 K 上满足 \(\int f d\mu = \int_{\mathrm{I}} f(j(x)) dx\) （对 \(f \in \mathcal{K}(\mathrm{X}; \mathbf{R})\) ）的测度，它是正的且总质量为 1；试证 \(\mathbf{b}_{\mu}\) 是 I 上的 Lebesgue 测度，且 \(\int g d\mu \neq g(\mathbf{b}_{\mu})\) 。

\begin{enumerate}
\def\labelenumi{\arabic{enumi})}
\setcounter{enumi}{1}
\tightlist
\item
  设 X 是紧空间，P 是 X 上取值于 \(]-\infty,+\infty]\) 的下半连续数值函数所成的集。假定 P 包含有限常数；对每个函数 \(h\in P\) 及 X 上每个正测度 \(\mu\) ，\(\mu^{*}(h)\) 有定义且 \(>-\infty\) （§4 习题 5）。
\end{enumerate}

\begin{enumerate}
\def\labelenumi{\alph{enumi})}
\item
  对两个正测度 \(\mu, \nu\) （\(X\) 上的），把 \(\mu \prec \nu\) 定义为：\(\mu^{*}(h) \leqslant \nu^{*}(h)\) 对每个函数 \(h \in \mathcal{P}\) 成立。试证关系 \(\mu \prec \nu\) 是 \(\mathcal{M}_{+}(X)\) 上的一个预序关系；它蕴涵 \(\mu(1) = \nu(1)\) 、\(c\mu \prec c\nu\) （对一切 \(c > 0\) ），以及 \(\mu + \lambda \prec \nu + \lambda\) 对每个测度 \(\lambda \in \mathcal{M}_{+}(X)\) 成立。
\item
  假定 \(\mathcal{P} + \mathcal{P} \subset \mathcal{P}\) ，且 \(c \cdot \mathcal{P} \subset \mathcal{P}\) 对 \(c > 0\) 成立。设 \(\mu\) 是 \(X\) 上一个正测度，\(f\) 是 \(\mathcal{C}(X; \mathbf{R})\) 中一个函数。用 \(Q_f\) 表示函数 \(h \geqslant f\) （\(\mathcal{P}\) 中的）所成的集，用 \(M_\mu\) 表示测度 \(\lambda \in \mathcal{M}_+(X)\) （满足 \(\lambda \prec \mu\) 者）所成的集。试证
\end{enumerate}

\[
\sup _ {\lambda \in \mathrm{M} _ {\mu}} \lambda (f) = \inf _ {h \in \mathrm{Q} _ {f}} \mu^ {*} (h).
\]

（对每个函数 \(g \in \mathcal{C}(\mathrm{X};\mathbf{R})\) ，令 \(p(g) = \inf_{h \in \mathrm{Q}_g} \mu^*(h)\) 。试证：\(p(g + g') \leqslant p(g) + p(g')\) 与 \(p(cg) = c \cdot p(g)\) 对 \(g, g'\) （\(\mathcal{C}(\mathrm{X};\mathbf{R})\) 中的）及 \(c > 0\) 成立；另一方面证明 \(\mathrm{M}_{\mu}\) 恰与满足下列条件的测度 \(\lambda \in \mathcal{M}_{+}(\mathrm{X})\) 所成之集相同：对每个函数 \(g \in \mathcal{C}(\mathrm{X};\mathbf{R})\) 有 \(\lambda(g) \leqslant p(g)\) ，然后应用 Hahn-Banach 定理得出结论。）若 \(\mathcal{P}\) 中函数连续且 \(S_0(f)\) 是 \(Q_f\) 的下包络，则也有 \(\mu(S_0(f)) = \sup \lambda(f)\) 。

\[
\lambda \in \mathrm{M} _ {\mu}
\]

\begin{enumerate}
\def\labelenumi{\alph{enumi})}
\setcounter{enumi}{2}
\item
  若集 \(\mathcal{P}_0\) （由 \(\mathcal{P}\) 中在 X 上连续且有限的函数组成）在 \(\mathcal{C}(\mathrm{X};\mathbf{R})\) 中完全，则关系 \(\mu \prec \nu\) 是 \(\mathcal{M}_{+}(\mathrm{X})\) 上的一个序关系。若测度 \(\nu \in \mathcal{M}_{+}(\mathrm{X})\) 对此序关系是极大的，则每个测度 \(\nu'\) （满足 \(0 \leqslant \nu' \leqslant \nu\) 者，对通常的序）也是极大的（用反证法并利用 \(a\) )）。若 \(\mathcal{P}_0 = \mathcal{P}\) 且 \(\mathcal{P}\) 完全，则关于序关系 \(\mu \prec \nu\) 的每个递增有向集对此关系在 \(\mathcal{M}_{+}(\mathrm{X})\) 中有上确界；特别地，\(\mathcal{M}_{+}(\mathrm{X})\) 对此关系是归纳的。
\item
  假定 \(\mathcal{P}\) 中函数连续且有限，\(\mathcal{P}\) 在 \(\mathcal{C}(\mathrm{X};\mathbf{R})\) 中完全，且 \(\mathcal{P} + \mathcal{P} \subset \mathcal{P}\) 与 \(c \cdot \mathcal{P} \subset \mathcal{P}\) 对一切 \(c > 0\) 成立。对每个函数 \(f \in \mathcal{C}(\mathrm{X};\mathbf{R})\) ，用 \(\mathrm{S}(f)\) 表示函数 \(h \in -\mathcal{P}\) （满足 \(f \leqslant h\) ）的下包络；对每个 \(\varepsilon > 0\) ，用 \(\mathrm{K}_{f,\varepsilon}\) 表示由 \(x \in \mathrm{X}\) （使 \(\mathrm{S}(f)(x) \geqslant f(x) + \varepsilon\) 者）组成的集。试证下列条件等价：
\end{enumerate}

\(\alpha)\) 测度 \(\mu \in \mathcal{M}_{+}(\mathrm{X})\) 关于序关系 \(\lambda \prec \lambda'\) 是极大的。

\(\beta)\) 对每个函数 \(f \in \mathcal{C}(\mathrm{X}; \mathbf{R})\) ，\(\mu(\mathrm{S}(f)) = \mu(f)\) 。

\(\beta^{\prime})\) 对每个函数 \(f \in \mathcal{P}\) ，\(\mu(\mathrm{S}(f)) = \mu(f)\) 。

\(\gamma)\) 对每个函数 \(f \in \mathcal{C}(\mathrm{X}; \mathbf{R})\) 及每个 \(\varepsilon > 0\) ，\(\mu(\mathrm{K}_{f,\varepsilon}) = 0\) 。

\(\gamma^{\prime})\) 对每个函数 \(f\in \mathcal{P}\) 及每个 \(\varepsilon >0\) ，\(\mu (\mathrm{K}_{f,\varepsilon}) = 0\)

（把 b) 应用于 \(-\mathcal{P}\) 。）

\begin{enumerate}
\def\labelenumi{\alph{enumi})}
\setcounter{enumi}{4}
\tightlist
\item
  在 \(d\) ) 的假设下，试证：若 \(\mu\) 与 \(\mu'\) 是 \(X\) 上对关系 \(\prec\) 极大的两个正测度，则 \(\mu + \mu'\) 亦然。由此推出：对此序关系极大的正测度所成之集对通常的序关系 \(\leqslant\) 是一个格。
\end{enumerate}

¶ 3) 设 X 是非空紧空间，P 是 X 上取值于 \([-\infty, +\infty]\) 的下半连续数值函数所成的集；假定 P 包含有限常数。称点 \(x \in X\) 为 P-极值点，如果测度 \(\varepsilon_{x}\) 关于预序关系 \(\mu \prec \nu (*)\) 是极小的；P-极值点所成之集记作 \(\mathrm{Ch}_{\mathcal{P}}(\mathrm{X})\) 。

\begin{enumerate}
\def\labelenumi{\alph{enumi})}
\tightlist
\item
  试证：若 \(\mathcal{P}'\) 是线性组合 \(\sum_{i} c_i h_i\) （\(\mathcal{P}\) 中函数的）所成之集
\end{enumerate}

（系数 \(\geqslant 0\) ）所成之集，则 \(\mathcal{P}\) -极值点与 \(\mathcal{P}'\) -极值点相同。

\begin{enumerate}
\def\labelenumi{\alph{enumi})}
\setcounter{enumi}{1}
\tightlist
\item
  对点 \(x \in X\) ，试证下列条件等价：
\end{enumerate}

\(\alpha)\) \(x\) 是 \(\mathcal{P}\) -极值点。

\(\beta)\) 对每个函数 \(f \in \mathcal{C}(\mathrm{X}; \mathbf{R})\) 及每个 \(\varepsilon > 0\) ，存在 \(h \in \mathcal{P}'\) 使 \(f \leqslant h\) 且 \(h(x) \leqslant f(x) + \varepsilon\) 。

\(\gamma)\) 对 \(x\) 的每个开邻域 U 及每个 \(\varepsilon > 0\) ，存在函数 \(h \geqslant 0\) 属于 \(\mathcal{P}'\) ，使 \(h(x) \leqslant \varepsilon\) 且 \(h(y) \geqslant 1\) 对一切 \(y \in X - U\) 成立。（利用习题 2 b。）

此外，试证 \(\mathrm{Ch}_{\mathcal{P}}(\mathrm{X})\) 中使 \(\mathcal{P}\) 至少有一个函数在 \(\mathrm{X}\) 达到其下确界的点所成之集在 \(\mathrm{Ch}_{\mathcal{P}}(\mathrm{X})\) 中稠密。

\begin{enumerate}
\def\labelenumi{\alph{enumi})}
\setcounter{enumi}{2}
\item
  称 X 的子集 F 为 P-稳定的，如果它是闭的，且关系 \(\lambda \prec \mu\) 与 \(\operatorname{Supp}(\mu) \subset F\) 蕴涵 \(\operatorname{Supp}(\lambda) \subset F\) 。试证：若 \(u \in P\) 且 \(F'\) 是 F 中使 u 在 F 达到其下确界的点所成之集，则 \(F'\) 是 P-稳定的。
\item
  再假定 \(\mathcal{P}\) 分隔 X 的点。试证：对每个函数 \(h \in \mathcal{P}\) ，集 \(S_h\) （X 中使 \(h\) 在 X 达到其下确界的点所成者）与 \(\mathrm{Ch}_{\mathcal{P}}(X)\) 相交。（考察一个极小 \(\mathcal{P}\) -稳定子集（含于 \(S_h\) 者），并用 \(c\) ) 证明这样的子集化为一点。）
\item
  在与 \(d\) ) 相同的假设下，试证：为使 \(\mathbf{F}\) （\(\mathbf{X}\) 的闭子集）包含 \(\mathrm{Ch}_{\mathcal{P}}(\mathbf{X})\) ，必须且只须对每个函数 \(h \in \mathcal{P}'\) ，\(\mathbf{F}\) 与 \(\mathbf{S}_h\) （使 \(h\) 在 \(\mathbf{X}\) 中达到其下确界的点所成之集）相交（利用 \(b\) )）。由此推出：为使 \(a \in \mathbf{X}\) 属于 \(\mathrm{Ch}_{\mathcal{P}}(\mathbf{X})\) 的闭包，必须且只须对每个开邻域 \(\mathbf{U}\) （含 \(a\) ）都存在 \(h \in \mathcal{P}'\) 与 \(b \in \mathbf{U}\) ，使 \(h(b) < h(x)\) 对一切 \(x \in \mathbf{X} - \mathbf{U}\) 成立（换言之，\(\mathbf{S}_h \subset \mathbf{U}\) ）。
\item
  在 \(d\) ) 的假设下，称 X 的子集 A 为 \(\mathcal{P}\) -边，如果对每个函数 \(h \in \mathcal{P}'\) ，\(\mathrm{S}_h\) 与 A 相交。取 X 为 \(\mathbf{R}^2\) 中以 \((-1,0)\) 与 \((1,0)\) 为心、以 1 为半径的两圆之并，取 \(\mathcal{P}\) 为 \(\mathbf{R}^2\) 上仿射线性函数在 X 上的限制所成之集。试证 \(\operatorname{Ch}_{\mathcal{P}}(\mathrm{X})\) 由 X 的点 \((\xi, \eta)\) （其中 \(|\xi| \geqslant 1\) ）组成。若 \(a\) 是 \(\mathcal{P}\) -极值点 \((1,1)\) ，试证存在 \(a\) 的一个邻域 U，使得没有函数 \(h \in \mathcal{P}'\) 满足 \(h \geqslant 0\) 、\(h(a) = 0\) 且 \(h(y) \geqslant 1\) 在 \(\mathrm{X} - \mathrm{U}\) 上成立。若 A 与 \(\mathrm{A}'\) 分别是 \(\operatorname{Ch}_{\mathcal{P}}(\mathrm{X})\) 中点 \((1,1)\) 与 \((-1,1)\) 的余集，则 A 与 \(\mathrm{A}'\) 都是 \(\mathcal{P}\) -边，但 \(\mathrm{A} \cap \mathrm{A}'\) 则不然；因而不存在最小的 \(\mathcal{P}\) -边。
\end{enumerate}

\(g)\) 设 \(\mathbf{Y}\) 是第二个紧空间，\(\mathcal{Q}\) 是 \(\mathbf{Y}\) 到 \([-\infty, +\infty]\) 中的下半连续映射所成之集，且包含有限常数。设 \(\mathcal{H}\) 为形如 \(h = f \otimes g\) （\(f \in \mathcal{P}\) ，\(g \in \mathcal{Q}\) ）的函数所成之集；试证 \(\mathrm{Ch}_{\mathcal{H}}(\mathrm{X} \times \mathrm{Y}) = \mathrm{Ch}_{\mathcal{P}}(\mathrm{X}) \times \mathrm{Ch}_{\mathcal{Q}}(\mathrm{Y})\) 。

¶ 4) 设 E 是 Hausdorff 局部凸空间，X 是 E 中的紧凸集，P 是在 X 上连续且凸的有限数值函数所成的集；\(\mu \prec \nu\) 表示由 P 在 \(\mathcal{M}_{+}(X)\) 上定义的预序关系（习题 2）。

\begin{enumerate}
\def\labelenumi{\alph{enumi})}
\item
  试证 \(\mu \prec \nu\) 是 \(\mathcal{M}_{+}(\mathrm{X})\) 上的一个序关系（利用 TVS, II, §5 习题 29），并且，对每个 \(x \in \mathrm{X}\) ，关系 \(\varepsilon_x \prec \mu\) 等价于说 \(\mu\) 的质量为 1 且 \(x\) 是 \(\mu\) 的重心。于是，对每个质量为 1 的（关于前述序关系的）极大测度 \(\mu\) ，存在一个且仅一个 \(x \in \mathrm{X}\) 使 \(\varepsilon_x \prec \mu\) 。反之，每个 \(z \in \mathrm{X}\) 都是至少一个极大测度的重心。为使 \(x \in \mathrm{X}\) 是 \(\mathrm{X}\) 的极值点，必须且只须 \(\varepsilon_x\) 是极大测度。
\item
  设 \(\mu\) 是 \(\mathrm{X}\) 上的极大测度，\(f\) 是 \(\mathcal{C}_{+}(\mathrm{X})\) 中的函数。试证：对每个 \(\varepsilon > 0\) ，存在函数 \(g\) 在 \(\mathrm{X}\) 上连续且凸，使 \(0 \leqslant g \leqslant f\) 且 \(\mu(g) \geqslant \mu(f) - \varepsilon\) （利用习题 3 b) 与 2 d)）。由此推出 \(\mu\) 的支集含于 \(\mathrm{X}\) 的极值点所成之集的闭包中。
\item
  设 \(\mu\) 是 \(X\) 上的极大测度，\((f_n)_{n \geqslant 1}\) 是 \(\mathcal{C}_+(X)\) 中函数的递减序列。试证：若对每个极值点 \(x \in X\) 序列 \((f_n(x))\) 趋于 0，则 \(\lim_{n \to \infty} \mu(f_n) = 0\) （利用 \(b\) ) 构造连续凸函数的递减序列 \((g_n)\) ，使 \(0 \leqslant g_n \leqslant f_n\) 与 \(\mu(g_n) \geqslant \mu(f_n) - \varepsilon\) 对一切 \(n\) 成立，并证明序列 \((g_n(y))\) 对每个 \(y \in X\) 趋于 0）。
\item
  由 c) 推出：若 \(A \subset X\) 不含任何极值点，且是 X 的紧子集序列 \((\mathrm{K}_{n})\) 之并，其中每个紧子集都是 X 中开集的可数交，则 \(\mu(\mathrm{A}) = 0\) 对每个极大测度 \(\mu\) 成立。
\item
  设 \(\mathcal{A}\) 是 X 上的连续函数所成的向量空间，这些函数是 E 上连续仿射线性函数在 X 上的限制。关系 \(\lambda \prec \mu\) 蕴含等式 \(\lambda(h) = \mu(h)\) 对每个函数 \(h \in \mathcal{A}\) 成立；对每个函数 \(f \in \mathcal{C}(\mathrm{X};\mathbf{R})\) ，\(\mathrm{S}(f)\) （函数 \(h \in \mathcal{A}\) 中满足 \(f \leqslant h\) 者的下包络）也是函数 \(g \in -\mathcal{P}\) 中满足 \(f \leqslant g\) 者的下包络。对每个 \(x \in X\) ，关系 \(\varepsilon_x \prec \mu\) 等价于 \(h(x) = \langle h, \mu \rangle\) 对一切 \(h \in \mathcal{A}\) 成立；X 的极值点与 \(\mathcal{A}\) -极值点相同。若 \(f \in \mathcal{P}\) ，则对每个 \(x \in X\) ，
\end{enumerate}

\[
\mathrm{S} (f) (x) = \sup _ {\varepsilon_ {x} \prec \mu} \int f (y) d \mu (y)\tag{*}
\]

（利用习题 2 b)）。若只假定 f 在 X 上凸且上半连续，则此公式不一定再成立（用习题 1 b) 的记号，考虑在 I 在 K 中的像上等于 1、而在其余处等于 0 的函数 f）。

\begin{enumerate}
\def\labelenumi{\alph{enumi})}
\setcounter{enumi}{5}
\tightlist
\item
  试证：当把右端限于以 x 为重心的离散测度 \(\mu\) 时，公式 (*) 仍然成立。
\end{enumerate}

¶ 5) 采用习题 4 e) 的记号，设 \(\mathcal{A}'\) 表示赋范空间 \(\mathcal{A}\) 的对偶，并赋以由 \(\mathcal{A}\) 的序结构导出的序结构（第二章，§2）。另一方面，用 C 表示由 \(X \times \{1\}\) 在空间 \(E \times R\) 中生成的以 0 为顶点的凸锥。

\begin{enumerate}
\def\labelenumi{\alph{enumi})}
\tightlist
\item
  试证下列性质等价：
\end{enumerate}

\begin{enumerate}
\def\labelenumi{(\roman{enumi})}
\item
  每个点 \(x \in X\) 都是唯一的、质量为 1 的极大测度 \(\beta_x\) 在 \(X\) 上的重心。
\item
  由 C 生成的向量空间 F，对于以 C 为其 \(\geqslant0\) 元素所成之集的序结构而言是一个格（换言之，X 是一个单形（TVS, II, §2 习题 41））。
\item
  对每个函数 \(f \in \mathcal{P}\) ，\(S(f)\) 既是凹的又是凸的。
\item
  若 \(\mu\) 与 \(\nu\) 是两个极大测度，且 \(\mu(h) = \nu(h)\) 对一切 \(h \in \mathcal{A}\) 成立，则 \(\mu = \nu\) 。
\item
  向量空间 \(\mathcal{A}'\) 是一个 Riesz 空间。
\item
  若 \(f\) 与 \(g\) 是 \(\mathcal{P}\) 中两个函数，则 \(S(f + g) = S(f) + S(g)\) 。
\end{enumerate}

（首先证明 (i) 蕴含下述性质：

\begin{enumerate}
\def\labelenumi{(\roman{enumi})}
\setcounter{enumi}{6}
\tightlist
\item
  映射 \(x \mapsto \beta_{x}\) 只以一种方式延拓为 F 到 \(\mathcal{M}(X)\) 中由极大测度生成的子空间上的双射线性映射，且此映射把 C 变换为极大测度所成的锥。
\end{enumerate}

然后利用习题 2 e) 由 (vii) 推出 (ii)。为从 (ii) 推出 (iii)，利用 Riesz 空间 F 中的分解引理及习题 4 f)。为从 (iii) 推出 (iv)，利用下述事实：若 \(g \in \mathcal{C}(X; \mathbf{R})\) 既凹又凸，则 \(h \in A\) 中使 \(h(x) > g(x)\) 对一切 \(x \in X\) 成立者所成之集是一个递减有向集（TVS, II, §5 命题 6），且以 g 为其下包络；然后应用习题 4 e)。为从 (iv) 推出 (v)，注意借助习题 4 e)，\(A'\) 可等同于 \(\mathcal{M}(X)\) 中由极大测度生成的子空间。为从 (v) 推出 (vi)，应用 Riesz 空间 \(A'\) 中的分解引理及习题 4 f)。最后，为从 (vi) 推出 (i)，考虑映射 \(f \mapsto S(f)(x)\) ，其中 \(f \in P\) ，\(x \in X\) ，并利用习题 4 e)。

\begin{enumerate}
\def\labelenumi{\alph{enumi})}
\setcounter{enumi}{1}
\item
  当 a) 的条件满足时，试证：对每个函数 \(f \in \mathcal{C}(X; \mathbf{R})\) ，映射 \(x \mapsto \beta_x(f)\) 属于 X 上两个上半连续函数之差所成的有界函数集在 X 上一致收敛拓扑下的闭包（利用 P 在 \(\mathcal{C}(X; \mathbf{R})\) 中完全这一事实）；因而该映射对 X 上每个测度可测。*对 X 上每个正测度 \(\mu\) ，唯一使 \(\nu \succ \mu\) 的极大测度由 \(\nu = \int \beta_x d\mu(x)\) 给出（参看第五章）。*
\item
  当 \(a\) ) 的条件满足且 X 此外还是可度量化时，试证：对每个 \(x \in X\) ，\(\beta_x\) 是唯一的正测度 \(\mu\) ，质量为 1、重心为 \(x\) ，且使 \(\mu(X - L) = 0\) ，其中 \(L\) 是极值点所成之集 \(\mathrm{Ch}_{\mathcal{A}}(X)\) （按定理 1 的方式论证）。
\item
  当 a) 的条件满足时，试证下列性质等价：
\end{enumerate}

\(\alpha)\) 集 \(\mathrm{L} = \mathrm{Ch}_{\mathcal{A}}(\mathrm{X})\) 是闭的，换言之 \(\check{\mathrm{S}}_{\mathcal{A}}(\mathrm{X}) = \mathrm{Ch}_{\mathcal{A}}(\mathrm{X})\) 。

\(\beta)\) 对每个函数 \(f \in \mathcal{P}\) ，\(\mathrm{S}(f)\) 在 \(\overline{\mathrm{L}}\) 上的限制连续。

\(\gamma)\) 对每个函数 \(f \in \mathcal{P}\) ，函数 \(S(f)\) 在 \(X\) 上连续。

\(\delta)\) 映射 \(x \mapsto \beta_x\) （X 到 \(\mathcal{M}_+(X)\) 中，赋以淡拓扑）连续。

（为证 \(\alpha\) ) 蕴含 \(\beta\) ) ，注意在 L 中 \(\beta_x = \varepsilon_x\) 。为证 \(\beta\) 蕴含 \(\gamma\) ) ，利用 S(f) 是 \(\mathcal{A}\) 中函数的递减有向集的下包络这一事实（TVS, II, §5 命题 6）、Dini 定理以及习题 3 \(d\) ) 应用于 \(\mathcal{A}\) 中的函数。为证 \(\gamma\) 蕴含 \(\delta\) ，利用习题 4 \(e\) )。最后，为证 \(\delta\) 蕴含 \(\alpha\) ，利用习题 4 \(a\) 。）

¶ 6) 设 X 是非空紧空间，H 是 \(\mathcal{C}(X;R)\) 的包含常数且分离 X 的点的线性子空间。于是，由 H 定义的（习题 2）关系 \(\lambda \prec \mu\) 是一个等价关系，且习题 3 意义下的 H-极值点与第 3 号意义下的 H-极值点相同。

设 \(\mathbf{F}\) 是 \(\mathbf{X}\) 的包含 \(\operatorname{Ch}_{\mathcal{H}}(\mathbf{X})\) 的闭子集；对每个 \(x \in \mathbf{X}\) ，用 \(\mathcal{M}_x^{\mathrm{F}}\) 表示正测度 \(\mu\) （在 \(\mathbf{X}\) 上、总质量为 1、使 \(\varepsilon_x \prec \mu\) 且 \(\operatorname{Supp}(\mu) \subset \mathbf{F}\) ）所成之集；为使 \(\varepsilon_x \in \mathcal{M}_x^{\mathrm{F}}\) ，必须且只须 \(x \in \mathbf{F}\) ；关系 \(x \in \operatorname{Ch}_{\mathcal{H}}(\mathbf{X})\) 等价于 \(\mathcal{M}_x^{\mathrm{F}} = \{\varepsilon_x\}\) 。对每个有界数值函数 \(f\) （定义在 \(\mathbf{F}\) 上）及每个 \(x \in X\) ，令

\[
\overline {{\mathrm{H}}} _ {x} ^ {\mathrm{F}} (f) = \inf _ {f \leqslant h | \mathrm{F}, h \in \mathscr {H}} h (x), \quad \underline {{\mathrm{H}}} _ {x} ^ {\mathrm{F}} (f) = - \overline {{\mathrm{H}}} _ {x} ^ {\mathrm{F}} (- f) = \sup _ {f \geqslant h | \mathrm{F}, h \in \mathscr {H}} h (x),
\]

它们都是有限数。

\begin{enumerate}
\def\labelenumi{\alph{enumi})}
\item
  试证：映射 \(f \mapsto \overline{\mathrm{H}}_x^{\mathrm{F}}(f)\) （\(\mathcal{C}(\mathrm{F};\mathbf{R})\) 到 \(\mathbf{R}\) 中的映射）是递增的、正齐次的且凸的。若 \(h \in \mathcal{H}\) ，则 \(\overline{\mathrm{H}}_x^{\mathrm{F}}(h) = h(x)\) 。
\item
  对每个 \(x \in X\) 、每个测度 \(\mu \in \mathcal{M}_x^{\mathrm{F}}\) 及每个函数 \(f \in \mathcal{C}(\mathrm{F};\mathbf{R})\) ，有 \(\underline{\mathrm{H}}_x^{\mathrm{F}}(f) \leqslant \int f d\mu \leqslant \overline{\mathrm{H}}_x^{\mathrm{F}}(f)\) 。反之，若对函数 \(f_0 \in \mathcal{C}(\mathrm{F};\mathbf{R})\) ，\(\gamma\) 是使 \(\underline{\mathrm{H}}_x^{\mathrm{F}}(f_0) \leqslant \gamma \leqslant \overline{\mathrm{H}}_x^{\mathrm{F}}(f_0)\) 的实数，则存在测度 \(\mu \in \mathcal{M}_x^{\mathrm{F}}\) 使 \(\gamma = \int f_0 d\mu\) 。
\item
  试证：对函数 \(g \in \mathcal{C}(\mathrm{X};\mathbf{R})\) ，下列条件等价：
\end{enumerate}

\(\alpha)\) 对每个 \(x \in X\) 及每个测度 \(\mu \in \mathcal{M}_x^{\mathrm{F}}\) ，有 \(g(x) = \int g d\mu\) 。

\(\beta)\) 对每个 \(x \in \mathrm{X}\) ，\(\underline{\mathrm{H}}_x^{\mathrm{F}}(g|\mathrm{F}) = \overline{\mathrm{H}}_x^{\mathrm{F}}(g|\mathrm{F}) = g(x)\) 。

\(\gamma)\) 对每个 \(\varepsilon > 0\) ，存在两个有限序列 \((h_i'), (h_j'')\) （\(\mathcal{H}\) 中的函数），使得令 \(h' = \sup(h_i')\) ，\(h'' = \inf(h_j'')\) ，就有 \(h' \leqslant g \leqslant h''\) 且 \(h'' - h' \leqslant \varepsilon\) 。

（为证 \(\beta\) ) 蕴含 \(\gamma\) ，按 GT, X, §4 第 1 号命题 2 的方式论证。）

当函数 \(g \in \mathcal{C}(\mathrm{X};\mathbf{R})\) 具有前述等价性质时，称它为 \(\mathcal{H}\) -调和的；\(\mathcal{H}^c\) （\(\mathcal{H}\) -调和函数所成之集）是 \(\mathcal{C}(\mathrm{X};\mathbf{R})\) 的包含 \(\mathcal{H}\) 的闭线性子空间；它与所考虑的闭集 \(\mathrm{F} \supset \mathrm{Ch}_{\mathcal{H}}(\mathrm{X})\) 无关。我们有 \((\mathcal{H}^c)^c = \mathcal{H}^c\) ，且 \(\mathcal{M}_+(\mathrm{X})\) 中由 \(\mathcal{H}^c\) 定义的等价关系与关系 \(\lambda \prec \mu\) （由 \(\mathcal{H}\) 所定义）相同；从而 \(\mathrm{Ch}_{\mathcal{H}^c}(\mathrm{X}) = \mathrm{Ch}_{\mathcal{H}}(\mathrm{X})\) 。映射 \(g \mapsto g|\mathrm{F}\) 是 \(\mathcal{H}^c\) 到它在 \(\mathcal{C}(\mathrm{F};\mathbf{R})\) 中的像上的严格递增等距（因而该像是 \(\mathcal{C}(\mathrm{F};\mathbf{R})\) 的闭子空间）。

\begin{enumerate}
\def\labelenumi{\alph{enumi})}
\setcounter{enumi}{3}
\item
  取 X 为 R 中区间 \([0,1]\) ，取 H 为 R 上二次多项式在 X 上的限制所成的向量空间。试证 \(H^{c}\) 不同于闭包 \(\overline{H}\) （H 在 \(\mathcal{C}(X;R)\) 中的闭包）。
\item
  试证：为使 \(\mathrm{Ch}_{\mathcal{H}}(\mathrm{X}) = \mathrm{X}\) ，必须且只须 \(\mathcal{H}^c = \mathcal{C}(\mathrm{X};\mathbf{R})\) （为证条件是必要的，利用 \(b\) )）。
\item
  试证：若 \(\mathcal{H}\) 是格序的（换言之，是一个 Riesz 空间），则 \(\mathcal{H}^c = \overline{\mathcal{H}}\) 。举出一个 \(\mathcal{H}\) 非格序而 \(\mathcal{H}^c = \overline{\mathcal{H}}\) 的例子。
\item
  设 \(\mathcal{E}_{\mathcal{H}}\) 是 \(\mathcal{C}(\mathrm{X};\mathbf{R})\) 中包含 \(\mathcal{H}\) 的最小闭子集，使 \(\inf (u,v)\) （\(\mathcal{E}_{\mathcal{H}}\) 中任意两个函数的下包络）属于 \(\mathcal{E}_{\mathcal{H}}\) 。试证下列性质等价：
\end{enumerate}

\(\alpha)\) \(f\in \mathcal{E}_{\mathcal{H}}\)

\(\beta)\) 对每个 \(x \in X\) 及每个测度 \(\mu \in \mathcal{M}_x^X\) ，\(\int f d\mu \leqslant f(x)\) ；

\(\gamma)\ \overline{\mathrm{H}}_{x}^{\mathrm{X}}(f) = f(x)\ \text{对于每个 } x \in \mathrm{X};\)

\(\delta)\) 对每个 \(\varepsilon > 0\) ，存在有限序列 \((h_i)\) （\(\mathcal{H}\) 中的函数），使 \(f \leqslant \inf(h_i) \leqslant f + \varepsilon\) 。

（为证 \(\beta\) ) 蕴含 \(\gamma\) ，利用 \(b\) 。）

由此推出 \(\mathcal{E}_{\mathcal{H}}\) 是一个尖凸锥，且 \(\mathcal{E}_{\mathcal{H}} \cap (-\mathcal{E}_{\mathcal{H}}) = \mathcal{H}^c\) 。试证 \(\mathcal{E}_{\mathcal{H}}\) 中每个函数在 \(X\) 中的下确界都在 \(\mathrm{Ch}_{\mathcal{H}}(X)\) 的至少一个点上达到。

¶ 7) 设 Y 是非空紧空间，R 是 \(\mathcal{C}(\mathrm{Y};\mathbf{R})\) 的一个线性子空间，它包含常数、分离 Y 的点且是格序的（换言之，是一个 Riesz 空间）。

\begin{enumerate}
\def\labelenumi{\alph{enumi})}
\tightlist
\item
  设 \(\mathcal{N}\) 是 \(\mathcal{R}\) 的一个极大孤立子空间（第二章，§1 习题 4）；试证存在一个且仅一个点 \(y_0\) ，使正线性形式 \(\varphi : f \mapsto f(y_0)\)
\end{enumerate}

在 \(\mathcal{R}\) 上是格的（第二章，§2 习题 5 b)），且 \(\mathcal{N} = \overline{\varphi}(0)\) 。（为看出 Y 中至少存在一点，使满足 \(f \geqslant 0\) 且属于 \(\mathcal{N}\) 的所有函数在该点为零，用反证法，证明在相反的情形下，顾及 Y 的紧性以及 R 中序关系的定义，将有 N = R。为证明 Y 中使 N 中所有函数都为零的点所成之集 \(Z(N)\) 化为单独一点，利用下述事实：N 是 R 中的一个超平面（第二章，§2 习题 5），以及 R 分离 Y 的点且包含常数。）

\begin{enumerate}
\def\labelenumi{\alph{enumi})}
\setcounter{enumi}{1}
\item
  在 \(a\) ) 的假设下，试证 \(y_0 \in \check{\mathrm{S}}_{\mathcal{R}}(\mathrm{Y})\) 。（令 \(\mathrm{S} = \check{\mathrm{S}}_{\mathcal{R}}(\mathrm{Y})\) ，注意 \(\mathcal{R}'\) （\(\mathcal{C}(\mathrm{S}; \mathbf{R})\) 中由限制到 \(\mathrm{S}\) 上的 \(\mathcal{R}\) 中函数所构成的子空间）作为序向量空间典范地同构于 \(\mathcal{R}\) （习题 6 c)），并且借助 \(f \mapsto f|\mathrm{S}\) 的逆同构，正线性形式 \(f \mapsto f(y_0)\) 给出 \(\mathcal{R}'\) 上的一个格的正线性形式；然后把 \(a\) ) 应用于 \(\mathcal{R}'\) 。）
\item
  对每个孤立子空间 \(\mathcal{N}_0 \neq \mathcal{R}\) ，试证存在极大孤立子空间 \(\mathcal{N} \supset \mathcal{N}_0\) 。（利用下述事实：若一个孤立子空间包含一个常数函数 \(\neq 0\) ，则它等于 \(\mathcal{R}\) 。）
\item
  设 \(\mathbf{Z}\) 是诸集 \(\mathbf{Z}(\mathcal{N})\) （每个都化为一点）之并，其中 \(\mathcal{N}\) 取遍 \(\mathcal{R}\) 的极大孤立子空间所成之集；试证 \(\mathbf{Z} = \check{\mathbf{S}}_{\mathcal{R}}(\mathbf{Y})\) 。（首先证明 \(\mathbf{Z}\) 是闭的，注意若 \(y \notin \mathbb{Z}\) ，则线性形式 \(f \mapsto f(y)\) 在 \(\mathcal{R}\) 上不是格的，从而（利用第二章，§2 习题 5 b)）线性形式 \(f \mapsto f(y')\) 对一切点 \(y'\) （充分接近 \(y\) 者）也不是格的。然后利用第 3 号命题 7，证明每个函数 \(f_0 \in \mathcal{R}\) 的下确界 \(\alpha\) （在 \(\mathbf{Y}\) 中取的）在 \(\mathbf{Z}\) 的至少一个点上达到；为此，令 \(\mathrm{M} = f_0^{-1}(\alpha)\) ，考虑子空间 \(\mathcal{N}_0\) ，它由函数 \(f \in \mathcal{R}\) （形如 \(f_1 - f_2\) ）组成。
\end{enumerate}

\(\mathrm{M} = \overline{f}_0(\alpha)\) ，考虑子空间 \(\mathcal{N}_0\) ，它由函数 \(f \in \mathcal{R}\) （形如 \(f_1 - f_2\) ，其中 \(f_1 \geqslant 0\) ，\(f_2 \geqslant 0\) ，且 \(f_1\) 与 \(f_2\) 在 \(\mathrm{M}\) 上为零）组成；证明 \(\mathcal{N}_0\) 是孤立的，并利用 \(c\) 。）

¶ 8) 一般记号与假设同习题 6，令 \(S = \check{S}_{\mathcal{H}}(X)\) 。称函数 \(f \in \mathcal{C}(S; \mathbf{R})\) 为 \(\mathcal{H}\) -预解的，如果存在函数 \(g \in \mathcal{H}^c\) 使 \(f = g|S\) （此函数届时是唯一的）。

\begin{enumerate}
\def\labelenumi{\alph{enumi})}
\tightlist
\item
  试证：关于函数 \(f \in \mathcal{C}(\mathbf{S};\mathbf{R})\) 的下列条件等价：
\end{enumerate}

\begin{enumerate}
\def\labelenumi{(\roman{enumi})}
\item
  \(f\) 是 \(\mathcal{H}\) -预解的。
\item
  对每个 \(x \in X\) ，\(\underline{\mathrm{H}}_x^{\mathrm{S}}(f) = \overline{\mathrm{H}}_x^{\mathrm{S}}(f)\) 。
\item
  对每个 \(x \in X\) 及每对测度 \(\mu_1, \mu_2\) （属于 \(\mathcal{M}_x^S\) ），有 \(\int f d\mu_1 = \int f d\mu_2\) 。
\end{enumerate}

（利用习题 6 b) 与 c)。）

\begin{enumerate}
\def\labelenumi{\alph{enumi})}
\setcounter{enumi}{1}
\tightlist
\item
  试证下列性质等价：
\end{enumerate}

\begin{enumerate}
\def\labelenumi{(\roman{enumi})}
\item
  每个函数 \(f \in \mathcal{C}(\mathrm{S};\mathbf{R})\) 都是 \(\mathcal{H}\) -预解的。
\item
  对每个 \(x \in X\) 及每个函数 \(f \in \mathcal{C}(\mathrm{S};\mathbf{R})\) ，\(\underline{\mathrm{H}}_x^{\mathrm{S}}(f) = \overline{\mathrm{H}}_x^{\mathrm{S}}(f)\) 。
\item
  集 \(\mathcal{M}_x^{\mathrm{S}}\) 化为单独一个元素。
\item
  集 \(\mathcal{H}^{\tilde{c}}\) 是格序的。
\item
  对每个函数 \(u \in \mathcal{E}_{\mathcal{H}}\) （习题 6 g)），存在最大下界 \(h_u\) （即 \(u\) 在 \(\mathcal{H}^c\) 中的最大下界）。
\end{enumerate}

（为证 (iv) 蕴含 (i)，对 \(\mathcal{H}^c\) 应用习题 7 d)，证明对每个点 \(s \in S\) ，线性形式 \(h \mapsto h(s)\) 在 \(\mathcal{H}^c\) 中是格的；然后利用 Stone 定理及习题 6 c)。为证 (i) 蕴含 (v)，注意 \(\mathcal{H}^c \subset \mathcal{E}_{\mathcal{H}}\) ，并考虑唯一的函数 \(h_u \in \mathcal{H}^c\) ，它使 \(h_u|S = u|S\) 。为证 (v) 蕴含 (iv)，注意若 \(h \in \mathcal{H}^c\) 则 \(-h^{-} = \inf(h,0) \in \mathcal{E}_{\mathcal{H}}\) ，并对 \(u = -h^{-}\) 应用 (v)。）

\begin{enumerate}
\def\labelenumi{\alph{enumi})}
\setcounter{enumi}{2}
\item
  若 \(b\) ) 的条件满足且 \(\gamma_x\) 是 \(\mathcal{M}_x^S\) 的唯一元素，试证映射 \(x \mapsto \gamma_x\) （\(X\) 到 \(\mathcal{M}_+(X)\) 中，赋以淡拓扑）连续；由此推出此时有 \(\mathrm{Ch}_{\mathcal{H}}(X) = \check{\mathrm{S}}_{\mathcal{H}}(X)\) （注意 \(\gamma_x = \varepsilon_x\) 在 \(\mathrm{Ch}_{\mathcal{H}}(X)\) 中成立）。
\item
  为使每个函数 \(f \in \mathcal{C}(\mathrm{S};\mathbf{R})\) 都是在 \(S\) 上的限制（某个函数 \(h \in \mathcal{H}\) 的限制），必须且只须 \(\mathcal{H}\) 在 \(\mathcal{C}(\mathrm{X};\mathbf{R})\) 中是闭的且是格序的。
\end{enumerate}

\begin{enumerate}
\def\labelenumi{\arabic{enumi})}
\setcounter{enumi}{8}
\tightlist
\item
  设 Y 是这样一个拓扑空间，其底集为乘积 \(I \times \{-1,0,1\}\) （在 \(R^{2}\) 中），其中 \(I = [0,1]\) ；对每个 \(a \in I\) 及每个 \(\varepsilon > 0\) ，令 \(U_{a,\varepsilon}\) 为由 \((x,y) \in Y\) （满足 \(|x - a| \leqslant \varepsilon\) 且 \((x,y)\) 异于 \((a,1)\) 与 \((a,-1)\) ）组成的集；诸集 \(\{(x,y)\}\) （\(y \neq 0\) 且 \(x \in I\) ）与诸集 \(U_{a,\varepsilon}\) 一起构成 Y 上一个拓扑的基，对此拓扑 Y 是一个不可度量化的紧空间（参看 GT, IX, §2 习题 13 d)）。用 X 表示 Y 与一个化为单点 \(\omega\) 的集的拓扑和空间。设 \(\mathcal{H}\) 为 X 上满足下式的连续有限数值函数 \(h\) 所成之集：
\end{enumerate}

\[
h (x, 0) = \frac {1}{2} \bigl (h (x, 1) + h (x, - 1) \bigr)
\]

对一切 \(x \in I\) 成立，且 \(h(\omega) = \int_0^1 h(x, 0) dx\) 。试证 \(\operatorname{Ch}_{\mathcal{H}}(X)\) 由点 \((x, y)\) （满足 \(y \neq 0\) 且 \(x \in I\) ）组成，但不存在 X 上质量为 1 的正测度 \(\mu\) 使 \(h(\omega) = \mu(h)\) 对一切 \(h \in \mathcal{H}\) 成立且 \(\mu(X - \operatorname{Ch}_{\mathcal{H}}(X)) = 0\) 。

¶ 10) a) 设 E 是 Hausdorff 完备局部凸空间，E′ 是其对偶，E′* ⊃ E 是 E′ 的代数对偶。设 A 是 E 的一个子集，它在弱化拓扑 σ(E, E′) 下是紧的。设 x 是 E′* 中的一点，属于 A 的凸包络 C 在弱拓扑 σ(E′,E′) 下的闭包。试证：若 \((x'_n)\) 是 E′ 中点的序列，在弱拓扑 σ(E′,E) 下收敛于 a′ ∈ E′，则 \(\lim_{n \to \infty}\langle x, x'_n\rangle = \langle x, a'\rangle\)。

（注意 \(x\) 是 A 上一个质量为 1 的测度的重心，并应用 Lebesgue 定理。）

\begin{enumerate}
\def\labelenumi{\alph{enumi})}
\setcounter{enumi}{1}
\item
  由 a) 推出：若 E 中存在一个在原拓扑下稠密的点序列，则限制到每个等度连续子集 \(H'\) （\(E'\) 的）上后，映射 \(x' \mapsto \langle x, x' \rangle\) 对弱拓扑 \(\sigma(\mathrm{E}', \mathrm{E})\) 连续（注意 \(H'\) 上由 \(\sigma(\mathrm{E}', \mathrm{E})\) 诱导的拓扑是可度量化的）。由此推出，此时必有 \(x \in E\) （TVS, IV, §5 第 5 号）；换言之，E 中一个在弱化拓扑下紧的集的闭凸包络对该拓扑仍然是紧的。
\item
  把 \(b\) ) 的结果推广到 \(E\) 是任意拟完备 Hausdorff 局部凸空间的情形（Krein 定理）。（首先化归为 \(E\) 是完备空间的情形，这通过考虑 \(\widehat{E}\) 实现；然后注意，凭 Eberlein 定理（TVS, IV, §5 第 3 号定理 1），只需证明每个序列 \((x_n)\) （\(C\) 中的点组成的）在 \(E\) 中对 \(\sigma(E, E')\) 有一个聚点；这就可化归为在 \(E\) 中存在一个原拓扑下稠密的序列的情形。）
\end{enumerate}

